\section{Extension: food-supply diversity after crop losses}\label{sec:divloss}
\paragraph{Purpose.} Section~\ref{sec:nutrition} asks how much of each nutrient a crop loss removes. This section asks
a different question: after the loss, does the country depend on \emph{fewer} or \emph{more evenly spread} food
sources, and how does that sit beside the change in total supply? The answer matters because a diversity index
can rise while nutrient availability falls.

\subsection{Method}
\paragraph{Categories.} Diversity is computed over the \emph{whole} national food supply: every FAOSTAT food-balance
(FBS) leaf item with positive supply, the same classification as the existing diversity table
(Section~\ref{sec:diversitymethod}) and as \citet{remans2014}. The item list is written to \file{food\_categories\_fbs.csv}
and is identical in the baseline and in every scenario (checked). The nine study crops are never renormalised
among themselves.

\paragraph{Supply after a loss.} Let $s_{i,n}$ be the 2021--23 national supply of nutrient $n$ from item $i$ (per
capita per day), and $f_c$ the share of crop $c$'s national supply removed (Eq.~\ref{eq:shock}: $f_c =
|\Delta Y_c|\,\text{SSR}^{*}_c$). An item of crop $c$ (primary product or its oil) keeps
\begin{equation}\label{eq:divitem}
s'_{i,n}=s_{i,n}\,(1-f_{c})\,q_{c,n}^{\,e_i},
\end{equation}
where $q_{c,n}$ is the elevated-CO$_2$ concentration factor of Section~\ref{sec:co2} ($q=1$ in scenarios A--C), and
$e_i=1$ for the primary product and $e_i=0$ for an extracted oil (no evidence transfers seed concentration changes
to oil). All other items are unchanged: no substitution, trade or household response is
modelled. Chickpea and lentil have no FBS item of their own; their supply $S_{c,n}$ (FAOSTAT SUA) is already inside
``Pulses, Other and products''. Their losses are therefore removed from that parent item and never added as
categories:
\begin{equation}\label{eq:divparent}
s'_{\text{parent},n}=s_{\text{parent},n}-\sum_{c\in\{\text{chickpea, lentil}\}} S_{c,n}\bigl[1-(1-f_c)\,q_{c,n}\bigr].
\end{equation}
Joint removals are checked against the parent supply; negative supply stops the run rather than being clipped.
In Belarus the SUA zinc and iron of chickpea plus lentil exceed the parent item by about 5\%, and in Canada the fat by
about 2\%, a mismatch between two sources compiled separately (the parent item is valued with one composition, the SUA
crops with their own; Section~\ref{sec:fbsmap}). No scenario comes close to removing the whole parent, so this does not bind, but it is
reported (check R4).

\paragraph{Diversity measures.} For each nutrient, with shares $p_{i,n}=s'_{i,n}/\sum_j s'_{j,n}$ over items with
positive supply:
\begin{equation}\label{eq:divH}
H_n=-\sum_i p_{i,n}\ln p_{i,n},\qquad N_{\text{eff},n}=e^{H_n},\qquad \text{top share}_n=\max_i p_{i,n}.
\end{equation}
Zero contributions are dropped ($0\ln0=0$). Total supply, NAR and MAR (Eqs.~\ref{eq:nar}) are reported beside them.

\paragraph{Why diversity and supply can move apart.} Differentiating $H$ with respect to one item's supply gives
\begin{equation}\label{eq:divrule}
\frac{\partial H}{\partial s_i}=-\frac{\ln p_i+H}{S},
\end{equation}
with $S$ the total. A small loss from item $i$ therefore \emph{raises} $H$ whenever $p_i>e^{-H}=1/N_{\text{eff}}$, that
is, whenever the item supplies more than an ``average'' source. Losing part of a dominant staple makes the
remaining shares more even while the total falls. In the single-crop scenarios where the crop is one FBS item, this
rule predicts the sign of the diversity change in every case (check R17, Appendix~\ref{app:ext}).

\paragraph{Scenarios.} All use the 2021--23 supply and the upper bound of Eq.~\ref{eq:shock} unless stated.
\begin{itemize}
\item \textbf{A. Single-crop standard loss:} one crop loses 10\% of domestic production (the existing shock; both
bounds). Scenario A reproduces the nutrient changes and MAR changes of Section~\ref{sec:nutrition} exactly (check R7).
\item \textbf{B. Simultaneous loss (stress test):} every modelled crop of a country loses 10\% at once. This is a
stress test, not a projection: such losses are not expected to coincide.
\item \textbf{C. Future climate, loss-only:} each modelled pair's conditional literature yield change
(Section~\ref{sec:futurenut}); gains are listed but not added, the convention of decision E14. Only the 59 modelled
pairs are shocked; nothing is extrapolated to other crops or countries. The yield change is measured from
1971--2000 to 2041--2060 and applied to 2021--23 supply, a \emph{hybrid counterfactual} (Section~\ref{sec:co2method}).
\item \textbf{C$'$. Present-relative variant:} the same with the change from the observed 2001--2024 level to
2041--2060.
\item \textbf{D. C combined with the elevated-CO$_2$ quality stress test} (Section~\ref{sec:co2}), with its coverage
limits.
\end{itemize}
Table~\ref{tab:divcover} shows which crops are modelled in each country, where the future scenario applies a loss and
which favourable changes are left out. Scenario C is therefore a \emph{loss-only stress test}, not a net projection:
yield gains are not added, and crops and countries outside the 59 pairs are not shocked.

\begin{table}[H]\centering\scriptsize
\caption[Crops covered by the diversity scenarios]{Crop--country pairs covered by the diversity scenarios. Scenario C
applies each pair's conditional yield loss; pairs with a projected gain are left out (loss-only).}\label{tab:divcover}
\adjustbox{max width=\linewidth}{\begin{tabular}{@{}lrp{4.6cm}p{3.6cm}p{3.0cm}@{}}
\toprule
Country & Pairs & Crops modelled & Loss applied in C & Gain left out of C \\
\midrule
Argentina & 1 & soybean & soybean & none \\
Australia & 2 & chickpea, lentil & chickpea, lentil & none \\
Bangladesh & 1 & rice & rice & none \\
Belarus & 1 & bean & bean & none \\
Brazil & 7 & rice, wheat, maize, soybean, bean, groundnut, rapeseed & bean, groundnut, maize, rapeseed, rice, soybean, wheat & none \\
Canada & 2 & lentil, rapeseed & lentil, rapeseed & none \\
China & 9 & rice, wheat, maize, soybean, chickpea, lentil, bean, groundnut, rapeseed & bean, chickpea, lentil, rapeseed, rice, wheat & maize, soybean, groundnut \\
Egypt & 1 & wheat & wheat & none \\
Ethiopia & 2 & chickpea, lentil & none & chickpea, lentil \\
France & 1 & rapeseed & rapeseed & none \\
India & 9 & rice, wheat, maize, soybean, chickpea, lentil, bean, groundnut, rapeseed & bean, chickpea, lentil, rapeseed, wheat & rice, maize, soybean, groundnut \\
Italy & 1 & rice & rice & none \\
Myanmar & 1 & bean & bean & none \\
Nigeria & 2 & rice, groundnut & none & rice, groundnut \\
Russia & 4 & wheat, soybean, chickpea, lentil & chickpea, lentil, wheat & soybean \\
South Africa & 3 & maize, soybean, rapeseed & rapeseed & maize, soybean \\
Tanzania & 1 & bean & bean & none \\
T\"urkiye & 1 & chickpea & chickpea & none \\
USA & 9 & rice, wheat, maize, soybean, chickpea, lentil, bean, groundnut, rapeseed & bean, chickpea, groundnut, lentil, maize, rapeseed, rice, soybean, wheat & none \\
Ukraine & 1 & maize & maize & none \\
\bottomrule
\end{tabular}
}
\end{table}

\subsection{Results}
\paragraph{India.} India's energy supply is spread over an effective \divIndBaseEffKcal\ items and its iron supply
over \divIndBaseEffIron, against \divWorldEffKcal\ and \divWorldEffIron\ for the world; the largest single sources
supply \divIndBaseTopKcal\% of energy (rice) and \divIndBaseTopIron\% of iron (wheat). A narrow base means that a
loss in one crop moves national supply and diversity visibly. Table~\ref{tab:divindia} shows how the scenarios change
both quantities.
\begin{itemize}
\item \textbf{Losing part of a dominant staple lowers supply but evens out the sources.} A 10\% wheat loss cuts
iron supply by \divIndWheatIronPctAbs\% yet raises the effective number of iron sources by \divIndWheatIronDEffAbs,
because wheat's iron share is far above $1/N_{\text{eff}}$. Read alone, the diversity index would suggest an
improvement.
\item \textbf{Losing a minor crop lowers both.} Chickpea or common bean losses reduce supply slightly and also
reduce diversity, because each supplies less than an average share.
\item \textbf{The simultaneous stress test} lowers energy by \divIndBKcalPctAbs\%, zinc by \divIndBZincPctAbs\% and iron by
\divIndBIronPctAbs\%, while the effective number of sources rises for every nutrient; MAR falls by \divIndBMARAbs\
points.
\item \textbf{Future climate (loss-only)} gives the same pattern at smaller size. The present-relative variant is
smaller still, because about \refShareOccurred\% of the modelled 1971--2000 to 2041--2060 yield change is already
contained in the observed 2001--2024 climate (Section~\ref{sec:co2method}).
\item \textbf{Adding the elevated-CO$_2$ quality stress test (D)} lowers zinc by \divIndDZincPctAbs\% and iron by
\divIndDIronPctAbs\% and MAR by \divIndDMARAbs\ points, and \emph{raises} the effective number of iron sources by
\divIndDIronDEffAbs. Wheat and rice, the largest sources, carry \cotwoIndWheatRiceMin--\cotwoIndWheatRiceMax\% of
the quality loss, so the remaining supply is more evenly spread while less nutritious. This is the clearest case in the results where the diversity index and
adequacy point in opposite directions.
\end{itemize}

\begin{table}[H]\centering\small
\caption[India: supply and diversity after losses]{India: change in national supply (\%) / change in effective
number of sources $N_{\text{eff}}$, and MAR change (points), by scenario. Upper bound of the supply pass-through.}
\label{tab:divindia}
\adjustbox{max width=\linewidth}{\begin{tabular}{@{}lrrrrr@{}}
\toprule
Scenario (India) & Energy & Protein & Zinc & Iron & $\Delta$MAR \\
\midrule
A: Wheat $-10\%$ & $-$2.11 / +0.41 & $-$2.38 / +0.34 & $-$3.26 / +0.52 & $-$2.95 / +0.58 & $-$0.52 \\
A: Rice $-10\%$ & $-$2.47 / +0.56 & $-$1.80 / +0.19 & $-$2.09 / +0.21 & $-$1.15 / +0.07 & $-$0.20 \\
A: Chickpea $-10\%$ & $-$0.23 / $-$0.02 & $-$0.51 / +0.01 & $-$0.34 / $-$0.01 & $-$0.40 / +0.02 & $-$0.07 \\
A: Common bean $-10\%$ & $-$0.13 / $-$0.03 & $-$0.31 / $-$0.04 & $-$0.20 / $-$0.03 & $-$0.44 / $-$0.03 & $-$0.08 \\
B: all nine $-10\%$ (stress test) & $-$5.74 / +0.79 & $-$5.46 / +0.45 & $-$6.40 / +0.64 & $-$5.50 / +0.58 & $-$0.98 \\
C: future climate, loss-only (hybrid) & $-$1.16 / +0.11 & $-$1.55 / +0.08 & $-$1.72 / +0.15 & $-$2.03 / +0.19 & $-$0.36 \\
C$'$: same, present-relative & $-$0.54 / +0.03 & $-$0.78 / +0.01 & $-$0.81 / +0.04 & $-$1.03 / +0.06 & $-$0.18 \\
D: C + elevated CO$_2$ quality & $-$1.16 / +0.11 & $-$4.45 / +0.44 & $-$5.44 / +0.68 & $-$4.20 / +0.49 & $-$0.75 \\
\bottomrule
\end{tabular}
}
\end{table}

\paragraph{Other countries.} Table~\ref{tab:divcountries} and Figure~\ref{fig:divsupply} compare the 20 panel
countries. Across all scenario, country and nutrient cases in which total supply falls, the effective number of
sources \emph{rises} in \divNUp\ of \divNCases\ (\divPctUp\%).
\begin{itemize}
\item \textbf{Where a modelled crop is the dominant source}, as rice is for energy in Bangladesh
(\divBdTopKcalShare\% of energy, \divBdEffKcal\ effective sources) and wheat is for iron in Russia
(\divRuTopIronShare\%), a loss raises diversity although few crops are shocked. In Bangladesh the stress test lowers
energy by \divBdBKcalPctAbs\% and raises the effective number of sources by \divBdBKcalDEffAbs; in Russia the future-climate
scenario (C) lowers iron by \divRuCIronPctAbs\% and raises the effective number of iron sources by \divRuCIronDEffAbs.
\item \textbf{Many staples shocked at once} gives the largest gains in effective sources: China (energy
\divCnBKcalPct\%, effective sources \divCnBKcalDEff) and India (\divIndBKcalPct\%, \divIndBKcalDEff).
\item \textbf{Diversity and supply can diverge across nutrients within one scenario.} Tanzania's bean loss under the
future climate lowers energy by \divTzCKcalPctAbs\% and the effective number of energy sources by a hair
(\divTzCKcalDEff), yet it lowers iron by \divTzCIronPctAbs\% while raising the effective number of iron sources by
\divTzCIronDEffAbs. Beans supply \divTzBeanIronShare\% of Tanzania's iron, above the threshold $1/N_{\text{eff}}=
\divTzThreshIron\%$, but only \divTzBeanKcalShare\% of its energy, below its threshold of \divTzThreshKcal\%
(Eq.~\ref{eq:divrule}).
\item \textbf{Where the lost crop is a minor source}, as for maize in Ukraine (iron \divUaCIronPct\%, effective
sources \divUaCIronDEff) or beans in Myanmar (iron \divMmCIronPct\%, \divMmCIronDEff), diversity and supply fall
together.
\item \textbf{Single-crop panels} (Argentina for soybean, France for rapeseed, Italy for rice) change little in
either measure, because the crop supplies only a small share of national energy (\divArgSoyKcalShare\%,
\divFrRapeKcalShare\% and \divItRiceKcalShare\%). The supply pass-through itself is close to full there
(SSR$^{*}\approx1$).
\end{itemize}

\begin{figure}[H]\centering
\includegraphics[width=\linewidth]{X04_diversity_vs_supply.pdf}
\caption[Diversity change against supply change]{Change in national supply (\%) against change in the effective number of
sources, for energy and iron, in scenarios B and C. Points above zero: the remaining sources are more even although
supply fell. Labels mark India and countries with losses above 3\%.}\label{fig:divsupply}
\end{figure}

\begin{table}[H]\centering\scriptsize
\caption[Supply and diversity after losses, by country]{Supply and diversity after losses by country: crops shocked,
baseline effective number of energy sources, and the changes in scenario B (stress test) and C (future climate,
loss-only, hybrid). $\Delta N$ = change in effective number of sources.}\label{tab:divcountries}
\adjustbox{max width=\linewidth}{\begin{tabular}{@{}lrrrrrrrrr@{}}
\toprule
Country & Crops & $N_{\text{eff}}$ energy & B energy \% & B $\Delta N$ & B iron \% & B $\Delta N$ & C iron \% & C $\Delta N$ & C $\Delta$MAR \\
\midrule
Argentina & 1 & 17.9 & $-$0.20 & $-$0.04 & 0.00 & 0.00 & 0.00 & 0.00 & 0.00 \\
Australia & 2 & 24.7 & $-$0.10 & $-$0.03 & $-$0.16 & $-$0.02 & $-$0.13 & $-$0.02 & 0.00 \\
Bangladesh & 1 & 5.9 & $-$6.29 & +0.53 & $-$3.62 & +0.54 & $-$0.02 & 0.00 & 0.00 \\
Belarus & 1 & 23.0 & 0.00 & 0.00 & 0.00 & 0.00 & 0.00 & 0.00 & 0.00 \\
Brazil & 7 & 22.8 & $-$3.57 & +0.49 & $-$4.37 & +0.58 & $-$4.00 & +0.46 & $-$0.78 \\
Canada & 2 & 25.4 & $-$1.02 & +0.20 & $-$0.11 & $-$0.03 & $-$0.03 & $-$0.01 & 0.00 \\
China & 9 & 19.1 & $-$4.61 & +0.99 & $-$3.16 & +0.18 & $-$0.21 & +0.01 & 0.00 \\
Egypt & 1 & 12.1 & $-$1.90 & +0.33 & $-$2.57 & +0.32 & $-$0.14 & +0.02 & 0.00 \\
Ethiopia & 2 & 14.0 & $-$0.15 & $-$0.01 & $-$0.22 & +0.01 & 0.00 & 0.00 & 0.00 \\
France & 1 & 23.8 & $-$0.20 & $-$0.04 & $-$0.01 & $-$0.01 & 0.00 & 0.00 & 0.00 \\
India & 9 & 16.1 & $-$5.74 & +0.79 & $-$5.50 & +0.58 & $-$2.03 & +0.19 & $-$0.36 \\
Italy & 1 & 20.8 & $-$0.16 & $-$0.04 & $-$0.07 & $-$0.02 & 0.00 & 0.00 & 0.00 \\
Myanmar & 1 & 8.0 & $-$0.09 & $-$0.02 & $-$0.34 & $-$0.04 & $-$0.74 & $-$0.09 & $-$0.13 \\
Nigeria & 2 & 19.6 & $-$1.38 & +0.01 & $-$0.88 & $-$0.04 & 0.00 & 0.00 & 0.00 \\
Russia & 4 & 16.8 & $-$3.05 & +0.76 & $-$4.94 & +0.86 & $-$3.61 & +0.62 & 0.00 \\
South Africa & 3 & 16.9 & $-$2.59 & +0.44 & $-$3.32 & +0.38 & 0.00 & 0.00 & 0.00 \\
T\"urkiye & 1 & 17.4 & $-$0.11 & $-$0.02 & $-$0.15 & $-$0.01 & $-$0.04 & 0.00 & 0.00 \\
Ukraine & 1 & 17.6 & $-$0.50 & $-$0.02 & $-$0.64 & $-$0.03 & $-$1.07 & $-$0.05 & 0.00 \\
Tanzania & 1 & 17.0 & $-$0.55 & $-$0.01 & $-$1.96 & +0.21 & $-$3.47 & +0.37 & $-$0.58 \\
USA & 9 & 22.6 & $-$3.67 & +0.65 & $-$3.99 & +0.84 & $-$2.61 & +0.48 & 0.00 \\
\bottomrule
\end{tabular}
}
\end{table}

\subsection{Reading diversity with adequacy}
Three readings are therefore not supported by these numbers:
\begin{enumerate}
\item \textbf{That every crop loss reduces diversity.} It does so only when the crop is a below-average source.
\item \textbf{That higher Shannon diversity is a nutritional improvement.} In the cases above it rises because the
largest source shrank; the total supply of every nutrient fell.
\item \textbf{That national food-source diversity describes diets.} $N_{\text{eff}}$ describes the composition of
national supply. It is not household dietary diversity, which is usually a count of foods or food groups eaten by
a household or person over a reference period \citep{ruel2003}, nor food preferences or eating behaviour.
\end{enumerate}
\citet{remans2014} used Shannon diversity of national food supply together with a functional measure based on the
nutrient composition of the items (closed access; not re-read for this extension). Functional diversity, developed
for farms by \citet{remans2011}, rises only when a food with a distinct nutrient profile is added; our item-level
Shannon measure does not distinguish nutritionally similar items, so it is not a measure of nutritional functional
diversity. \citet{khoury2014} found that over 50 years national food supplies became more even in their crop
commodities while becoming more similar to each other worldwide, so a rising within-country evenness can coexist
with a narrowing of the global crop base. Our results add a short-term counterpart: within a country, a loss in the
dominant staple raises evenness while lowering supply. For adequacy, total supply and NAR/MAR are the decisive
quantities; diversity describes how exposed the supply is to any single source.

\paragraph{Limitations.} No substitution, imports or price response; national means hide distribution within
countries; MAR is capped at 1 per nutrient, so it moves only where a nutrient is near or below its requirement;
FBS items differ in breadth (``Vegetables, other'' is one item), so $N_{\text{eff}}$ depends on the classification.
